\section{A finite quadratic sieve with explicit coefficients}
\label{sec:sieve}

Fix \(z\ge1\), and put
\[
 \sfD_z=\{d:1\le d\le z,\ d\text{ squarefree}\}.
\]
This set contains \(1\) and is closed under taking positive divisors.
For any squarefree \(d\), define the positive rationals
\begin{equation}
 g(d)=\frac{\rho(d)}d,\qquad
 h(d)=\prod_{p\mid d}\frac{\rho(p)}{p-\rho(p)}.
 \label{eq:gh}
\end{equation}
The prime denominators are positive: \(2-\rho(2)=1\), and
\(p-\rho(p)=p-2\ge1\) for odd primes.
By the empty-product convention \(g(1)=h(1)=1\).
Set
\[
 G=G(z)=\sum_{r\in\sfD_z}h(r)\ge1,\qquad
 W_d=\sum_{\substack{r\in\sfD_z\\d\mid r}}h(r).
\]
For \(d\in\sfD_z\), choose
\begin{equation}
 \boxed{\qquad
 \lambda_d=\frac{\mu(d)W_d}{g(d)G}.
 \qquad}
 \label{eq:weights}
\end{equation}
These are explicit finite rational expressions. There is no minimization
over real vectors in their definition.

\begin{lemma}[Normalization and coefficient bound]
\label{lem:weights}
The coefficients satisfy
\begin{equation}
 \lambda_1=1,\qquad
 |\lambda_d|g(d)=\frac{W_d}G\le1,\qquad
 |\lambda_d|\rho(d)\le d.
 \label{eq:coefficient-bound}
\end{equation}
For every \(r\in\sfD_z\),
\begin{equation}
 \sum_{\substack{d\in\sfD_z\\r\mid d}}\lambda_dg(d)
   =\frac{\mu(r)h(r)}G.
 \label{eq:transform}
\end{equation}
\end{lemma}

\begin{proof}
The normalization follows from \(W_1=G\).
Since \(d\) is squarefree, \(|\mu(d)|=1\), and \(0<W_d\le G\);
this proves the bounds. For the transform, expand \(W_d\) and change the
order of two finite sums:
\begin{align*}
 \sum_{\substack{d\in\sfD_z\\r\mid d}}\lambda_dg(d)
 &=\frac1G
   \sum_{\substack{v\in\sfD_z\\r\mid v}}
          h(v)\sum_{r\mid d\mid v}\mu(d)\\
 &=\frac{\mu(r)h(r)}G.
\end{align*}
The last equality is Lemma~\ref{lem:mobius}. All divisors occurring in the
inner sum belong to \(\sfD_z\), because \(v\le z\).
\end{proof}

\begin{lemma}[Quadratic identity]
\label{lem:quadratic}
The coefficients in \eqref{eq:weights} satisfy
\begin{equation}
 Q:=\sum_{d,e\in\sfD_z}\lambda_d\lambda_e
             g(\lcm(d,e))=\frac1{G(z)}.
 \label{eq:quadratic}
\end{equation}
\end{lemma}

\begin{proof}
For a prime \(p\), the definition of \(h\) gives
\[
 1+\frac1{h(p)}=\frac1{g(p)}.
\]
Expanding a finite product over the primes dividing a squarefree integer
\(v\) yields
\[
 \frac1{g(v)}
   =\prod_{p\mid v}\left(1+\frac1{h(p)}\right)
   =\sum_{r\mid v}\frac1{h(r)}.
\]
Multiplicativity on squarefree prime supports also gives
\[
 g(\lcm(d,e))=\frac{g(d)g(e)}{g(\gcd(d,e))}
   =g(d)g(e)\sum_{r\mid\gcd(d,e)}\frac1{h(r)}.
\]
Insert this into \(Q\) and interchange finite sums:
\begin{align*}
 Q
 &=\sum_{r\in\sfD_z}\frac1{h(r)}
       \left(\sum_{\substack{d\in\sfD_z\\r\mid d}}
                     \lambda_dg(d)\right)^2\\
 &=\sum_{r\in\sfD_z}\frac1{h(r)}
                 \left(\frac{\mu(r)h(r)}G\right)^2
   =\frac1{G^2}\sum_{r\in\sfD_z}h(r)
   =\frac1G.
\end{align*}
Here \(\mu(r)^2=1\) on \(\sfD_z\), and the transform is
\eqref{eq:transform}. This is an equality of finite rational expressions.
\end{proof}

\subsection{The sieve majorant and its error}

Let \(S(X,z)\) count integers \(1\le n\le X\) for which no prime at most
\(z\) divides \(n(n+2)\). This is a bounded condition.
For such an \(n\), the only member \(d\in\sfD_z\) dividing \(n(n+2)\)
is \(1\): any \(d>1\) has a prime divisor at most \(d\le z\).
Thus the following pointwise inequality holds for every \(n\):
\begin{equation}
 \mathbf 1_{\{n\text{ survives}\}}
 \le
 \left(\sum_{\substack{d\in\sfD_z\\d\mid n(n+2)}}\lambda_d\right)^2.
 \label{eq:majorant}
\end{equation}
If \(n\) survives, the right side is \(\lambda_1^2=1\); otherwise the
left side is \(0\) and a rational square is nonnegative.

Sum over \(1\le n\le X\), expand the square, and use the equivalence
\[
 d\mid n(n+2)\ \wedge\ e\mid n(n+2)
 \quad\Longleftrightarrow\quad
 \lcm(d,e)\mid n(n+2).
\]
We obtain
\[
 S(X,z)\le
 \sum_{d,e\in\sfD_z}\lambda_d\lambda_e A_{\lcm(d,e)}(X).
\]
By Lemma~\ref{lem:count-error}, write
\[
 A_v(X)=Xg(v)+E_v,\qquad |E_v|\le\rho(v)
\]
for the squarefree numbers \(v=\lcm(d,e)\).
There is no restriction \(v\le z\) here; indeed \(v\le z^2\).
Lemma~\ref{lem:quadratic} gives the main term \(X/G(z)\).
For the error, \eqref{eq:rho-lcm} and
\eqref{eq:coefficient-bound} give
\begin{align}
 \left|\sum_{d,e\in\sfD_z}\lambda_d\lambda_e E_{\lcm(d,e)}\right|
 &\le
 \sum_{d,e\in\sfD_z}|\lambda_d||\lambda_e|
                   \rho(d)\rho(e)\notag\\
 &=\left(\sum_{d\in\sfD_z}|\lambda_d|\rho(d)\right)^2
 \le\left(\sum_{d=1}^z d\right)^2
 \le z^4.
 \label{eq:sieve-error}
\end{align}
The final inequality uses \(z\) terms, each at most \(z\).
Consequently,
\begin{equation}
 S(X,z)\le\frac{X}{G(z)}+z^4.
 \label{eq:sifted}
\end{equation}

\begin{proposition}[Finite twin-prime bound]
\label{prop:finite-sieve}
For every \(X\ge0\) and \(z\ge1\),
\begin{equation}
 \boxed{\qquad
 T(X)\le z+\frac{X}{G(z)}+z^4.
 \qquad}
 \label{eq:finite-sieve}
\end{equation}
\end{proposition}

\begin{proof}
At most \(z\) pairs counted by \(T(X)\) have their smaller prime at most
\(z\). Every remaining pair has both primes greater than \(z\).
No prime at most \(z\) can divide either of those primes, and therefore
none divides their product. Its smaller member is counted by \(S(X,z)\).
Thus \(T(X)\le z+S(X,z)\), and \eqref{eq:sifted} proves the result.
\end{proof}
